\chapter{Introdução à semântica operacional}

\section{Introdução}

Até agora, estudamos como reconhecer se uma sequência de tokens
pertence a uma linguagem: a análise léxica e sintática. O próximo
passo é atribuir \textbf{significado} aos programas reconhecidos:
dado um programa sintaticamente correto, o que ele \emph{faz}? Essa
pergunta é o objeto de estudo da \textbf{semântica de linguagens de
programação}.

Existem várias maneiras de formalizar o significado de um programa.
Neste capítulo, estudaremos a \textbf{semântica operacional}, que
define o significado de um programa descrevendo como ele é
\emph{executado}, passo a passo, por uma máquina abstrata. A
semântica operacional é particularmente útil para a construção de
compiladores e interpretadores, pois a sua especificação se traduz
diretamente em código.

Apresentamos dois estilos de semântica operacional, ambos para uma
linguagem de expressões aritméticas com adição, multiplicação e
números inteiros:

\begin{itemize}
  \item
    \textbf{Small-step} (\emph{pequenos passos}): define a execução
    como uma sequência de reduções elementares, cada uma
    simplificando a expressão um pouco.
  \item
    \textbf{Big-step} (\emph{grandes passos}): define diretamente o
    valor final de uma expressão, abstraindo os passos intermediários.
\end{itemize}

Por fim, mostramos como a semântica big-step se traduz em um
interpretador Haskell, disponível em

\begin{flushleft}
  \verb|Exp/Interp/ExpInterp.hs|
\end{flushleft}

\section{A Linguagem de Expressões}

A linguagem que estudaremos é formada por três construções: números
inteiros, adição e multiplicação. Sua sintaxe abstrata pode ser
descrita pela gramática:

\[
  \begin{array}{lcl}
    e & ::= & n \mid e_1 + e_2 \mid e_1 * e_2
  \end{array}
\]

onde \(n\) representa um número inteiro. Em Haskell, o tipo de dado
correspondente é:

\begin{minted}{haskell}
data Exp
  = EInt Int
  | Exp :+: Exp
  | Exp :*: Exp
\end{minted}

Os construtores infixos \haskell{(:+:)} e
\haskell{(:*:)} refletem diretamente a estrutura
binária das operações.

\section{Semântica Small-Step}

A semântica \textbf{small-step} (ou semântica de redução) define a
execução de uma expressão como uma sequência de passos de
simplificação. Cada passo é descrito por uma \textbf{relação de
redução} \(e \to e'\), que se lê ``a expressão \(e\) reduz em um
passo para \(e'\)''.

\subsection{Valores}

Antes de apresentar as regras, precisamos distinguir as expressões
que já chegaram ao seu resultado final, os \textbf{valores},
das que ainda podem ser reduzidas. Para a linguagem de expressões
aritméticas, o único tipo de valor é um número inteiro:

\[
  v ::= n
\]

\subsection{Regras de Redução}

As regras de redução são apresentadas na forma de um sistema de
inferência. Cada regra tem a forma:

\[
  \infer[_{\{Nome\}}]{\text{conclusão}}{\text{premissas}}
\]

Se não houver premissas, a regra é um \textbf{axioma} e é escrita sem
a linha horizontal. Para a linguagem de expressões, temos as seguintes regras:

\textbf{Adição de valores:}

\[
  \infer[_{\{Add\}}]{n_1 + n_2 \to n_1 \mathbin{\hat{+}} n_2}{}
\]

onde \(\mathbin{\hat{+}}\) denota a adição aritmética de inteiros
(para distinguir o operador da linguagem do operador matemático).

\textbf{Redução no lado esquerdo da adição:}

\[
  \infer[_{\{AddL\}}]{e_1 + e_2 \to e_1' + e_2}{e_1 \to e_1'}
\]

\textbf{Redução no lado direito da adição:}

\[
  \infer[_{\{AddR\}}]{n_1 + e_2 \to n_1 + e_2'}{e_2 \to e_2'}
\]

A regra (AddR) só se aplica quando o lado esquerdo já é um valor
\(n_1\), garantindo uma ordem de avaliação da esquerda para a direita.

As regras para multiplicação são análogas:

\[
  \infer[_{\{Mul\}}]{n_1 * n_2 \to n_1 \mathbin{\hat{*}} n_2}{}
\]

\[
  \infer[_{\{MulL\}}]{e_1 * e_2 \to e_1' * e_2}{e_1 \to e_1'}
\]

\[
  \infer[_{\{MulR\}}]{n_1 * e_2 \to n_1 * e_2'}{e_2 \to e_2'}
\]

\subsection{Sequências de Redução}

A execução completa de uma expressão é uma sequência de reduções que
termina em um valor:

\[
  e \to e_1 \to e_2 \to \cdots \to v
\]

Escrevemos \(e \to^* v\) quando existe tal sequência (zero ou mais
passos). Para uma expressão que já é um valor, temos \(v \to^* v\)
com zero passos.

\subsection{Exemplo: Small-Step}

Considere a expressão \((2 + 3) * 4\). A redução small-step procede assim:

\[
  (2 + 3) * 4 \xrightarrow{\text{MulL}} 5 * 4 \xrightarrow{\text{Mul}} 20
\]

O passo 1 usa (MulL) com premissa \(2 + 3 \to 5\) (pela regra Add). O
passo 2 usa (Mul) diretamente, pois ambos os operandos já são valores.

Considere agora \(1 + 2 * 3\). Aplicando (AddL) e depois (MulL):

\[
  1 + 2 * 3
  \xrightarrow{\text{AddR}} 1 + 6
  \xrightarrow{\text{Add}} 7
\]

Note que a regra (AddR) se aplica porque o lado esquerdo \(1\) já é
um valor, forçando a avaliação do subtermos \(2 * 3\) antes da soma.

\section{Semântica Big-Step}

A semântica \textbf{big-step} (ou semântica natural) define
diretamente a relação entre uma expressão e o seu valor final,
abstraindo todos os passos intermediários. A relação é escrita \(e
\Downarrow v\), que se lê ``a expressão \(e\) avalia para o valor \(v\)''.

\subsection{Regras de Avaliação}

\textbf{Um número avalia para si mesmo:}

\[
  \infer[_{\{Num\}}]{n \Downarrow n}{}
\]

\textbf{Adição:}

\[
  \infer[_{\{Add\}}]{e_1 + e_2 \Downarrow n_1 \mathbin{\hat{+}}
  n_2}{e_1 \Downarrow n_1 & e_2 \Downarrow n_2}
\]

\textbf{Multiplicação:}

\[
  \infer[_{\{Mul\}}]{e_1 * e_2 \Downarrow n_1 \mathbin{\hat{*}}
  n_2}{e_1 \Downarrow n_1 & e_2 \Downarrow n_2}
\]

A estrutura é mais direta que a semântica small-step: cada regra
corresponde exatamente a um construtor da árvore sintática, e as
premissas avaliam recursivamente os sub-termos.

\subsection{Exemplos: Big-Step}

\textbf{Exemplo 1:} \((2 + 3) * 4 \Downarrow 20\)

A derivação é uma árvore de inferência:

\[
  \dfrac{
    \dfrac{
      \dfrac{}{2 \Downarrow 2} \quad \dfrac{}{3 \Downarrow 3}
    }{2 + 3 \Downarrow 5}
    \quad
    \dfrac{}{4 \Downarrow 4}
  }{(2 + 3) * 4 \Downarrow 20}
\]

\textbf{Exemplo 2:} \(1 + 2 * 3 \Downarrow 7\)

\[
  \dfrac{
    \dfrac{}{1 \Downarrow 1}
    \quad
    \dfrac{
      \dfrac{}{2 \Downarrow 2} \quad \dfrac{}{3 \Downarrow 3}
    }{2 * 3 \Downarrow 6}
  }{1 + 2 * 3 \Downarrow 7}
\]

Note que a árvore reflete diretamente a estrutura da árvore sintática
abstrata da expressão, é exatamente essa correspondência que torna
a semântica big-step tão conveniente para a implementação de interpretadores.

\section{Implementação em Haskell}

A semântica big-step se traduz em Haskell de forma quase literal.
Cada regra de inferência corresponde a uma equação da função
\haskell{eval}:

\begin{minted}{haskell}
module Exp.Interp.ExpInterp (eval) where

import Exp.Frontend.Syntax.ExpSyntax

eval :: Exp -> Int
eval (EInt n)    = n
eval (e1 :+: e2) = eval e1 + eval e2
eval (e1 :*: e2) = eval e1 * eval e2
\end{minted}

A correspondência com as regras semânticas é direta:

\begin{itemize}
  \item A equação \haskell{eval (EInt n) = n} implementa
    a regra (Num): um número \(n\) avalia para si mesmo.
  \item A equação \haskell{eval (e1 :+: e2) = eval e1 + eval e2} implementa
    a regra (Add): avalia recursivamente \(e_1\)
    e \(e_2\) e soma os resultados.
  \item A equação \haskell{eval (e1 :*: e2) = eval e1 * eval e2} implementa
    a regra (Mul) da mesma forma.
\end{itemize}

O casamento de padrões (pattern matching) de Haskell corresponde às
premissas de cada regra, e a chamada recursiva a
\haskell{eval} corresponde a aplicar a regra à sub-expressão.

\subsection{Executando o Interpretador}

O interpretador é parte do compilador
\haskell{exp}. Para usá-lo em modo de interpretação,
basta invocar:

\begin{verbatim}
cabal run exp -- --interp --file input.exp
\end{verbatim}

onde \verb|input.exp| contém uma expressão na
linguagem. Por exemplo, para a entrada

\begin{verbatim}
(2 + 3) * 4
\end{verbatim}

o interpretador produz \haskell{20}, confirmando o
resultado que calculamos pelas derivações semânticas acima.

\subsection{Totalidade e Terminação}

A função \haskell{eval} é \textbf{total}: está
definida para todas as expressões do tipo
\haskell{Exp} e sempre termina. Isso pode ser
verificado observando que cada chamada recursiva é feita sobre
sub-expressões estritamente menores, a estrutura da árvore
sintática abstrata diminui a cada nível. Essa propriedade é garantida
pela indução estrutural e corresponde, na semântica, ao fato de que
toda expressão da linguagem possui um valor (não há expressões que
ficam presas sem produzir resultado).

\section{Conclusão}

Neste capítulo, introduzimos a semântica operacional como ferramenta
para definir o significado de programas. Estudamos dois estilos
complementares:

\begin{itemize}
  \item
    A \textbf{semântica small-step} expõe cada passo de redução
    individualmente, permitindo raciocinar sobre o processo de execução.
  \item
    A \textbf{semântica big-step} relaciona expressões diretamente
    aos seus valores finais e se traduz quase literalmente em uma
    função recursiva Haskell.
\end{itemize}

A implementação do interpretador \haskell{eval}
demonstra que uma especificação semântica formal não é apenas um
documento teórico: ela é um projeto de implementação preciso. Nos
próximos capítulos, estenderemos esta linguagem com variáveis,
comandos e funções, exigindo que a semântica modele também
\textbf{ambientes} e \textbf{efeitos}.
